\section*{Chapitre \ref{ch:mx:build-a-matching-engine-and-exchange-simulator} --- Construction~: un moteur d'appariement et un simulateur de bourse}

\begin{solution}{exo:mx:build-a-matching-engine-and-exchange-simulator:1}
5\,013~: le paquet transporte les messages 5\,001 à 5\,012, et les heartbeats portent le prochain numéro attendu sans l'avancer.
\end{solution}

\begin{solution}{exo:mx:build-a-matching-engine-and-exchange-simulator:2}
Ligne A 5.3~\% ($308/5\,793$), ligne B 8.9~\% ($517/5\,793$), les deux 0.52~\% ($30/5\,793$).
\end{solution}

\begin{solution}{exo:mx:build-a-matching-engine-and-exchange-simulator:3}
L'ordre avait déjà été exécuté (à 5 secondes), il n'y avait donc rien à annuler~: motif «~L~», annulation trop tardive. Une annulation d'un ordre inconnu reçoit le
même motif.
\end{solution}

\begin{solution}{exo:mx:build-a-matching-engine-and-exchange-simulator:4}
Avec un seul générateur suivant l'ordre des événements, ajouter un participant ou modifier le tempo de l'un d'eux décale chaque tirage suivant, si bien que les latences de tous les autres participants et les
pertes du flux changent~: les comparaisons entre simulations cessent d'être comparables. Des flux issus de générateurs à compteur, indexés par clé, donnent à chaque usage sa propre suite.
\end{solution}

\begin{solution}{exo:mx:build-a-matching-engine-and-exchange-simulator:5}
Pour savoir quels messages incrémentaux l'instantané contient déjà~: appliquer ceux qui suivent son numéro de séquence et ignorer ceux qui le précèdent. Sans cela, un récepteur
applique certains messages deux fois ou en oublie, et le carnet reconstruit est faux d'une manière qui peut ne pas apparaître en haut du carnet.
\end{solution}

\begin{solution}{exo:mx:build-a-matching-engine-and-exchange-simulator:6}
Que les trois moteurs s'accordent sur chaque message, rapport et instantané produit à partir des entrées du jeu de test, qui exercent chaque type de message et chaque contrôle. Cela ne
prouve pas que les règles sont correctes (les trois pourraient partager une même mauvaise lecture d'une règle), ni qu'ils s'accordent sur des entrées absentes du jeu de test.
\end{solution}

\begin{solution}{exo:mx:build-a-matching-engine-and-exchange-simulator:7}
La ligne A ne perd rien, la ligne B perd les 166 messages de sa coupure, et aucune \omterm{def:mx:build-a-matching-engine-and-exchange-simulator:feed}{demande de retransmission} n'est nécessaire~: la ligne A livre chaque message.
\end{solution}

\begin{solution}{exo:mx:build-a-matching-engine-and-exchange-simulator:8}
Le temps du simulateur va aux agents, à la boucle d'événements, aux tâches de fond et aux paquets du flux, et non au seul moteur~; le chiffre C++ inclut le
démarrage du programme et ses vérifications, c'est donc une borne inférieure pour le seul moteur. L'accélération d'une session dépend de la part de son temps consacrée au moteur.
\end{solution}

\begin{solution}{pb:mx:build-a-matching-engine-and-exchange-simulator:1}
\textbf{1.} Voir les définitions.
\textbf{2.} Flux 14, saisie des ordres 5, \omterm{def:mx:build-a-matching-engine-and-exchange-simulator:oe}{rapports d'exécution} 6, contrôles 10 (auxquels s'ajoutent les messages de session propres à SoupBinTCP).
\textbf{3.} Acceptation~; une transaction d'auto-croisement~; un IOC avec 200 exécutés et 300 annulés~; un ordre post-only au repos~; une exécution au prix de l'ordre au repos~; un
rejet hors \omterm{def:mx:the-limit-order-book:tick}{pas de cotation}~; une annulation trop tardive~; une expiration de fin de journée.
\textbf{4.} Voir les définitions.
\textbf{5.} Priorité prix-temps~: dans l'ordre temporel au \omterm{def:mx:the-limit-order-book:tick}{niveau de prix}~; au prorata~: l'allocation paramétrable du livre~1, avec une part attribuée au premier ordre.
\textbf{6.} La limite de l'ordre, la bande de prix, les quantités minimales, la \omterm{def:mx:build-a-matching-engine-and-exchange-simulator:stp}{prévention de l'auto-appariement}.
\textbf{7.} Pré-ouverture et enchère d'ouverture, négociation en continu ou enchères par lots, enchère de clôture, suspendu, en pause (\omterm{def:mx:circuit-breakers-and-halts:pause}{enchère de réouverture})~; quitter une enchère revient à la dénouer.
\textbf{8.} Par les contrôles R (bande) et P (phase), déclenchés aux heures prévues.
\textbf{9.} Session, premier numéro de séquence, nombre, puis les messages~; A et B transportent les mêmes paquets.
\textbf{10.} À partir de flux issus de générateurs à compteur indexés par ligne~: pertes indépendantes, rafales de Gilbert--Elliott, doublons, gigue, coupures.
\textbf{11.} La ligne A a perdu 308 messages sur 5\,793, la ligne B 517~; 30 manquaient sur les deux lignes~; 30 \omterm{def:mx:build-a-matching-engine-and-exchange-simulator:feed}{demandes de retransmission} les ont comblés.
\textbf{12.} Les ordres d'un instantané de mi-journée, puis les messages incrémentaux postérieurs à son numéro de séquence~; comparaison avec le carnet du moteur sur dix niveaux de profondeur.
\textbf{13.} Traiter les deux flux, retenir chaque numéro de séquence dans le premier à l'arriver, et combler les trous sur les deux.
\textbf{14.} Une horloge simulée, un départage des égalités fixe, un aléa à compteur indexé par finalité, aucune horloge murale dans le moteur.
\textbf{15.} Quatre minutes de flux de transactions plus un agent chaos envoyant chaque type de message et chaque contrôle~: 17\,059 enregistrements de journal.
\textbf{16.} \emph{Résultat nommé}~: le moteur Python passe les treize scénarios de règles du moteur et reproduit le SHA-256 du jeu de test~; les moteurs C++20 et Rust
reproduisent octet pour octet la sortie du moteur Python sur les 17\,059 enregistrements du jeu de test (et font l'aller-retour sur 4\,866 messages de référence)~; sur ce portable, le moteur Python
traite environ 60\,000 enregistrements par seconde et le programme de test C++20 plus de 500\,000.
\textbf{17.} Elle inclut le démarrage du programme, le décodage et les comparaisons octet par octet.
\textbf{18.} Les agents, les modèles de latence, les dégradations du flux et le temps de réponse des serveurs en production.
\textbf{19.} Les spreads implicites, une passerelle FIX, les règles propres à chaque plateforme vérifiées dans la documentation de chacune.
\textbf{20.} Afin qu'une différence entre deux simulations ne puisse provenir que d'une différence dans leurs entrées.
\end{solution}

\begin{solution}{iq:mx:build-a-matching-engine-and-exchange-simulator:1}
Arbitrer les lignes A et B selon le numéro de séquence~; en cas de trou sur les deux, mettre en tampon les messages suivants et demander les manquants au service de retransmission~; si
le trou est trop ancien, reconstruire à partir d'un instantané et appliquer les messages postérieurs à son numéro de séquence.
\iqlookfor{Arbitrage, retransmission, instantanés.}
\end{solution}

\begin{solution}{iq:mx:build-a-matching-engine-and-exchange-simulator:2}
Une table associant chaque prix à son niveau (une structure triée ou un tableau indexé par ticks), chaque niveau étant une file FIFO d'ordres, et une table de hachage de la référence de l'ordre vers son nœud
afin qu'une annulation soit en O(1)~; garder les meilleurs prix en cache.
\iqlookfor{Niveaux, FIFO, référence d'ordre.}
\end{solution}

\begin{solution}{iq:mx:build-a-matching-engine-and-exchange-simulator:3}
Une boucle à événements discrets sur une horloge simulée, avec un départage déterministe des égalités, des flux aléatoires à compteur indexés par finalité, sans dépendance à l'horloge murale ni
à l'ordre des tables de hachage, et un journal de chaque entrée~; tester en hachant les sorties de simulations répétées.
\iqlookfor{Horloge, aléa, ordonnancement, tests.}
\end{solution}

\begin{solution}{iq:mx:build-a-matching-engine-and-exchange-simulator:4}
Une acceptation, un \omterm{def:mx:build-a-matching-engine-and-exchange-simulator:oe}{rapport d'exécution} par exécution (avec la quantité restante), puis une annulation du reliquat avec le motif IOC~; le flux montre les exécutions contre
les ordres au repos.
\iqlookfor{Séquence des rapports~; quantité restante.}
\end{solution}

\begin{solution}{iq:mx:build-a-matching-engine-and-exchange-simulator:5}
L'\omterm{def:mx:build-a-matching-engine-and-exchange-simulator:stp}{annulation à la déconnexion} empêche les ordres périmés de s'exécuter tant qu'un participant ne peut pas les gérer~; la \omterm{def:mx:build-a-matching-engine-and-exchange-simulator:stp}{prévention de l'auto-appariement} empêche les wash trades entre les stratégies
propres d'une même firme, ainsi que les frais et les problèmes de reporting qu'ils entraînent.
\iqlookfor{Ordres périmés~; wash trades.}
\end{solution}

\begin{solution}{iq:mx:build-a-matching-engine-and-exchange-simulator:6}
Rejouer le journal jusqu'au premier enregistrement dont la sortie diffère, décoder les deux sorties, relire la règle dans la spécification, et écrire un scénario minimal qui
le reproduit~; l'ajouter au jeu de test pour que les trois moteurs y soient soumis.
\iqlookfor{Dichotomie sur le journal~; un nouveau test.}
\end{solution}
